\Subhead{6.3.1}{Минимизация дисперсии ошибки. Принцип ортогональности ошибки и наблюдения}

Пусть $x,y$ — центрированные совместно стационарные процессы, а $\hat x(t)=\int_{-\infty}^{\infty}h(u)y(t-u)\dd u$. Минимизируем $J=\langle[x(t)-\hat x(t)]^2\rangle$. Вариация по $h$ даёт
\begin{equation}
 \langle e(t)y(t-\tau)\rangle=0,\quad e=x-\hat x,
\end{equation}
то есть ошибка ортогональна всему доступному линейному пространству наблюдений. Если бы корреляция с некоторым наблюдением оставалась, добавление его малого вклада могло бы уменьшить ошибку.

Определим $C_{xy}(\tau)=\langle x(t+\tau)y(t)\rangle$. Тогда нормальные уравнения
\begin{equation}
 \int h(u)C_y(\tau-u)\dd u=C_{xy}(\tau)
\end{equation}
в некаузальном случае дают после преобразования Фурье
\begin{equation}
 H_{\rm nc}(\omega)=\frac{S_{xy}(\omega)}{S_y(\omega)}.
\end{equation}
Для $y=x+n$, некоррелированных $x,n$, имеем $H=S_x/(S_x+S_n)$: частоты с большим отношением сигнал/шум пропускаются лучше. Это оптимальность по среднеквадратичной ошибке \emph{случайного сообщения}, а не максимум выходного отношения сигнал/шум для одного известного импульса. Минимальная ошибка $J_{\min}=C_x(0)-(2\pi)^{-1}\int|S_{xy}|^2/S_y\dd\omega$.
